% chapterintro
% ==============================
% 第0章 导言：0.0-chinese.md
学习CMake之前，需要对系统进行设置，这样才能运行所有示例。

本章的主要内容有：

\begin{itemize}
\item 如何获取代码

\item 如何在GNU/Linux、macOS和Windows上安装运行示例所需的所有工具

\item 自动化测试如何工作

\item 如何报告问题，并提出改进建议

\end{itemize}

我们会尽可能让初学者看懂本书的内容。不过，这本书并非完全适合零基础人士。我们假设，您对构建目标平台上可用的软件，及本地工具有基本的了解。有Git版本控制的经验，可与源码库进行“互动"(不是必需)。


% --- recipe 1: 0.1-chinese.md ---
\section{0.1 获取代码}

本书的源代码可以在GitHub上找到，网址是 \texthref{https://github.com/dev-cafe/cmake-cookbook} 。开源代码遵循MIT许可：只要原始版权和许可声明包含在软件/源代码的任何副本中，可以以任何方式重用和重新混合代码。许可的全文可以在 \texthref{https://opensource.org/licenses/MIT} 中看到。

为了测试源码，需要使用Git获取代码：

\begin{itemize}
\item 主要的GNU/Linux发行版都可以通过包管理器安装Git。也可以从Git项目网站 \texthref{https://git-scm.com} 下载二进制发行版，进行安装。

\item MacOS上，可以使用自制或MacPorts安装Git。

\item Windows上，可以从Git项目网站( \texthref{https://git-scm.com} )下载Git可执行安装文件。

\end{itemize}

可以通过GitHub桌面客户端访问这些示例，网址为 \texthref{https://desktop.github.com} 。

另一种选择是从 \texthref{https://github.com/dev-cafe/cmake-cookbook} 下载zip文件。

安装Git后，可以将远程库克隆到本地计算机，如下所示：

\begin{codebox}[shell]
$ git clone https://github.com/dev-cafe/cmake-cookbook.git
\end{codebox}

这将创建一个名为\texttt{cmake-cookbook}的文件夹。本书内容与源码的章节对应，书中章节的编号和源码的顺序相同。

在GNU/Linux、MacOS和Windows上，使用最新的持续集成进行测试。我们会在之后讨论测试的设置。

我们用标签v1.0标记了与本书中打印的示例相对应的版本。为了与书中内容对应，可以如下获取此特定版本：

\begin{codebox}[shell]
$ git clone --single-branch -b v1.0 https://github.com/dev-cafe/cmake-cookbook.git
\end{codebox}

我们希望收到Bug修复，并且GitHub库将继续发展。要获取更新，可以选择库的master分支。


% --- recipe 2: 0.2-chinese.md ---
\section{0.2 Docker镜像}

在Docker中进行环境搭建，无疑是非常方便的(依赖项都已经安装好了)。我们的Docker镜像是基于Ubuntu 18.04的镜像制作，您可以按照官方文档\texthref{https://docs.docker.com} 在您的操作系统上安装Docker。

Docker安装好后，您可以下载并运行我们的镜像，然后可以对本书示例进行测试:

\begin{codebox}[shell]
$ docker run -it devcafe/cmake-cookbook_ubuntu-18.04
$ git clone https://github.com/dev-cafe/cmake-cookbook.git
$ cd cmake-cookbook
$ pipenv install --three
$ pipenv run python testing/collect_tests.py 'chapter-*/recipe-*'
\end{codebox}


% --- recipe 3: 0.3-chinese.md ---
\section{0.3 安装必要的软件}

与在Docker中使用不同，另一种选择是直接在主机操作系统上安装依赖项。为此，我们概括了一个工具栈，可以作为示例的基础。您必须安装以下组件：

\begin{enumerate}
\item CMake

\item 编译器

\item 自动化构建工具

\item Python

\end{enumerate}

我们还会详细介绍，如何安装所需的某些依赖项。

\mySmallsection{0.3.1 获取CMake}

本书要使用的CMake最低需要为3.5。只有少数示例，演示了3.5版之后引入的新功能。每个示例都有提示，指出示例代码在哪里可用，以及所需的CMake的最低版本。提示信息如下:

\begin{myInfoBox}
\textbf{NOTE}：这个示例的代码可以在 \texthref{https://github.com/dev-cafe/cmake-cookbook/tree/v1.0/chapter-03/recipe10} 中找到，其中包括一个C示例。该示例在CMake 3.5版(或更高版本)中是有效的，并且已经在GNU/Linux、macOS和Windows上进行了测试。
\end{myInfoBox}

有些(如果不是大多数)示例仍然适用于较低版本的CMake。但是，我们没有测试过这个。我们认为CMake 3.5是大多数系统和发行版的默认软件，而且升级CMake也没什么难度。

CMake可以以多种方式安装。下载并提取由Kitware维护的二进制发行版，可以在所有平台上运行，下载页面位于 \texthref{https://cmake.org/download/} 。

大多数GNU/Linux发行版都在包管理器中提供了CMake。然而，在一些发行版中，版本可能比较旧，因此下载由Kitware提供的二进制文件当然是首选。下面的命令将从CMake打包的版本中下载并安装在\texttt{\$HOME/Deps/CMake}(根据您的偏好调整此路径)下的CMake 3.5.2：

\begin{codebox}[shell]
$ cmake_version="3.5.2"
$ target_path=$HOME/Deps/cmake/${cmake_version}
$ cmake_url="https://cmake.org/files/v${cmake_version%.*}/cmake-${cmake_version}-Linux-x86_64.tar.gz"
$ mkdir -p "${target_path}"
$ curl -Ls "${cmake_url}" | tar -xz -C "${target_path}" --strip-components=1
$ export PATH=$HOME/Deps/cmake/${cmake_version}/bin${PATH:+:$PATH}
$ cmake --version
\end{codebox}

macOS获取最新版本的CMake：

\begin{codebox}[shell]
$ brew upgrade cmake
\end{codebox}

Windows上，可以使用Visual Studio 2017，它提供了CMake支持。Visual Studio 2017的安装记录在第13章，\emph{可选生成器和交叉编译}，示例技巧1，\emph{使用Visual Studio 2017构建CMake项目}。

或者，可以从 \texthref{https://www.msys2.org} 下载MSYS2安装程序，按照其中给出的说明更新包列表，然后使用包管理器\texttt{pacman}安装CMake。下面的代码正在构建64位版本：

\begin{codebox}[shell]
$ pacman -S mingw64/mingw-w64-x86_64-cmake
\end{codebox}

对于32位版本，请使用以下代码(为了简单起见，我们以后只会提到64位版本)：

\begin{codebox}[shell]
$ pacman -S mingw64/mingw-w64-i686-cmake
\end{codebox}

MSYS2的另一个特性是在Windows上提供了一个终端，比较像Unix操作系统上的终端，提供可用的开发环境。

\mySmallsection{0.3.2 编译器}

我们将需要C++、C和Fortran的编译器。编译器的版本需要比较新，因为我们需要在大多数示例中支持最新的语言标准。CMake为来自商业和非商业供应商的许多编译器，提供了非常好的支持。为了让示例始终能够跨平台，并尽可能独立于操作系统，我们使用了开源编译器:

\begin{itemize}
\item GNU/Linux上，GNU编译器集合(GCC)是直接的选择。它是免费的，适用于所有发行版。例如，在Ubuntu上，可以安装以下编译器：

\begin{codebox}[shell]
$ sudo apt-get install g++ gcc gfortran
\end{codebox}

\item 在LLVM家族中，Clang也是C++和C编译器的一个很好的选择：

\begin{codebox}[shell]
$ sudo apt-get install clang clang++ gfortran
\end{codebox}

\item macOS上，XCode附带的LLVM编译器适用于C++和C。我们在macOS测试中使用了GCC的Fortran编译器。GCC编译器必须使用包管理器单独安装：

\begin{codebox}[shell]
$ brew install gcc
\end{codebox}

\item Windows上，可以使用Visual Studio测试C++和C示例。或者，可以使用MSYS2安装程序，MSYS2环境中(对于64位版本)使用以下单个命令安装整个工具链，包括C++、C和Fortran编译器：

\begin{codebox}[shell]
$ pacman -S mingw64/mingw-w64-x86_64-toolchain
\end{codebox}

\end{itemize}

\mySmallsection{0.3.3 自动化构建工具}

自动化构建工具为示例中的项目提供构建和链接的基础设施，最终会安装和使用什么，很大程度上取决于操作系统：

\begin{itemize}
\item GNU/Linux上，GNU Make(很可能)在安装编译器时自动安装。

\item macOS上，XCode将提供GNU Make。

\item Windows上，Visual Studio提供了完整的基础设施。MSYS2环境中，GNU Make作为mingw64/mingw-w64-x86\_64工具链包的一部分，进行安装。

\end{itemize}

为了获得最大的可移植性，我们尽可能使示例不受这些系统相关细节的影响。这种方法的优点是配置、构建和链接，是每个编译器的\emph{固有特性}。

Ninja是一个不错的自动化构建工具，适用于GNU/Linux、macOS和Windows。Ninja注重速度，特别是增量重构。为GNU/Linux、macOS和Windows预先打包的二进制文件可以在GitHub库中找到，网址是 \texthref{https://github.com/ninja-build/ninja/releases} 。

Fortran项目中使用CMake和Ninja需要注意。使用CMake 3.7.2或更高版本是必要的，Kitware还有维护Ninja，相关包可以在 \texthref{https://github.com/Kitware/ninja/releases} 上找到。

在GNU/Linux上，可以使用以下一系列命令安装Ninja：

\begin{codebox}[shell]
$ mkdir -p ninja
$ ninja_url="https://github.com/Kitware/ninja/releases/download/v1.8.2.g3bbbe.kitware.dyndep-1.jobserver-1/ninja-1.8.2.g3bbbe.kitware.dyndep-1.jobserver-1_x86_64-linux-gnu.tar.gz"
$ curl -Ls ${ninja_url} | tar -xz -C ninja --strip-components=1
$ export PATH=$HOME/Deps/ninja${PATH:+:$PATH}
\end{codebox}

Windows上，使用MSYS2环境(假设是64位版本)执行以下命令：

\begin{codebox}[shell]
$ pacman -S mingw64/mingw-w64-x86_64-ninja
\end{codebox}

\begin{myInfoBox}
\textbf{NOTE}：我们建议阅读这篇文章 \texthref{http://www.aosabook.org/en/posa/ninja.html} ，里面是对NInja编译器的历史和设计的选择，进行启发性的讨论。
\end{myInfoBox}

\mySmallsection{0.3.4 Python}

本书主要关于CMake，但是其中的一些方法，需要使用Python。因此，也需要对Python进行安装：解释器、头文件和库。Python 2.7的生命周期结束于2020年，因此我们将使用Python 3.5。

在Ubuntu 14.04 LTS上(这是Travis CI使用的环境，我们后面会讨论)，Python 3.5可以安装如下：

\begin{codebox}[shell]
sudo apt-get install python3.5-dev
\end{codebox}

Windows可使用MSYS2环境，Python安装方法如下(假设是64位版本):

\begin{codebox}[shell]
$ pacman -S mingw64/mingw-w64-x86_64-python3
$ pacman -S mingw64/mingw-w64-x86_64-python3-pip
$ python3 -m pip install pipenv
\end{codebox}

为了运行已经写好的测试机制，还需要一些特定的Python模块。可以使用包管理器在系统范围内安装这些包，也可以在隔离的环境中安装。建议采用后一种方法：

\begin{itemize}
\item 可以在不影响系统环境的情况下，将安装包进行清理/安装。

\item 可以在没有管理员权限的情况下安装包。

\item 可以降低软件版本和依赖项冲突的风险。

\item 为了复现性，可以更好地控制包的依赖性。

\end{itemize}

为此，我们准备了一个\texttt{Pipfile}。结合\texttt{pipfile.lock}，可以使用\texttt{Pipenv}( \texthref{http://pipenv.readthedocs} )。创建一个独立的环境，并安装所有包。要为示例库创建此环境，可在库的顶层目录中运行以下命令：

\begin{codebox}[shell]
$ pip install --user pip pipenv --upgrade
$ pipenv install --python python3.5
\end{codebox}

执行\texttt{pipenv shell}命令会进入一个命令行环境，其中包含特定版本的Python和可用的包。执行\texttt{exit}将退出当前环境。当然，还可以使用\texttt{pipenv run}在隔离的环境中直接执行命令。

或者，可以将库中的\texttt{requirements.txt}文件与\texttt{Virtualenv}( \texthref{http://docs.pythonguide.org/en/latest/dev/virtualenvs/} )和\texttt{pip}结合使用，以达到相同的效果：

\begin{codebox}[shell]
$ virtualenv --python=python3.5 venv
$ source venv/bin/activate
$ pip install -r requirements.txt
\end{codebox}

可以使用\texttt{deactivate}命令退出虚拟环境。

另一种选择是使用\texttt{Conda}环境，我们建议安装\texttt{Miniconda}。将把最新的\texttt{Miniconda}安装到GNU/Linux的\texttt{\$HOME/Deps/conda}目录(从 \texthref{https://repo.continuum.io/miniconda/miniconda3-latestlinux-x86\_64.sh} 下载)或macOS(从 \texthref{https://repo.continuum.io/miniconda/miniconda3-latestmacosx-x86\_64.sh} 下载)：

\begin{codebox}[shell]
$ curl -Ls https://repo.continuum.io/miniconda/Miniconda3-latest-Linux-x86_64.sh > miniconda.sh
$ bash miniconda.sh -b -p "$HOME"/Deps/conda &> /dev/null
$ touch "$HOME"/Deps/conda/conda-meta/pinned
$ export PATH=$HOME/Deps/conda/bin${PATH:+:$PATH}
$ conda config --set show_channel_urls True
$ conda config --set changeps1 no
$ conda update --all
$ conda clean -tipy
\end{codebox}

Windows上，可以从 \texthref{https://repo.continuum.io/miniconda/Miniconda3-latest-Windows-x86\_64.exe} 下载最新的\texttt{Miniconda}。该软件包可以使用\texttt{PowerShell}安装，如下:

\begin{codebox}[shell]
$basedir = $pwd.Path + "\"
$filepath = $basedir + "Miniconda3-latest-Windows-x86_64.exe"
$Anaconda_loc = "C:\Deps\conda"
$args = "/InstallationType=JustMe /AddToPath=0 /RegisterPython=0 /S /D=$Anaconda_loc"
Start-Process -FilePath $filepath -ArgumentList $args -Wait -Passthru
$conda_path = $Anaconda_loc + "\Scripts\conda.exe"
$args = "config --set show_channel_urls True"
Start-Process -FilePath "$conda_path" -ArgumentList $args -Wait -Passthru
$args = "config --set changeps1 no"
Start-Process -FilePath "$conda_path" -ArgumentList $args -Wait -Passthru
$args = "update --all"
Start-Process -FilePath "$conda_path" -ArgumentList $args -Wait -Passthru
$args = "clean -tipy"
Start-Process -FilePath "$conda_path" -ArgumentList $args -Wait -Passthru
\end{codebox}

安装了\texttt{Conda}后, Python模块可以按如下方式安装:

\begin{codebox}[shell]
$ conda create -n cmake-cookbook python=3.5
$ conda activate cmake-cookbook
$ conda install --file requirements.txt
\end{codebox}

执行\texttt{conda deactivate}将退出\texttt{conda}的环境。

\mySmallsection{0.3.5 依赖软件}

有些示例需要额外的依赖，这些软件将在这里介绍。

\mySmallsection{0.3.5.1 BLAS和LAPACK}

大多数Linux发行版都为BLAS和LAPACK提供包。例如，在Ubuntu 14.04 LTS上，您可以运行以下命令：

\begin{codebox}[shell]
$ sudo apt-get install libatlas-dev liblapack-dev liblapacke-dev
\end{codebox}

macOS上，XCode附带的加速库可以满足我们的需要。

Windows使用MSYS2环境，可以按如下方式安装这些库(假设是64位版本)：

\begin{codebox}[shell]
$ pacman -S mingw64/mingw-w64-x86_64-openblas
\end{codebox}

或者，可以从GitHub ( \texthref{https://github.com/referlapack/lapack} )下载BLAS和LAPACK的参考实现，并从源代码编译库。商业供应商为平台提供安装程序，安装包中有BLAS和LAPACK相关的API。

\mySmallsection{0.3.5.2 消息传递接口(MPI)}

MPI有许多商业和非商业实现。这里，安装免费的非商业实现就足够了。在Ubuntu 14.04 LTS上，我们推荐\texttt{OpenMPI}。可使用以下命令安装：

\begin{codebox}[shell]
$ sudo apt-get install openmpi-bin libopenmpi-dev
\end{codebox}

在macOS上，\texttt{Homebrew}发布了\texttt{MPICH}：

\begin{codebox}[shell]
$ brew install mpich
\end{codebox}

还可以从 \texthref{https://www.open-mpi.org/software/} 上获取源代码，编译\texttt{OpenMPI}。 对于Windows，Microsoft MPI可以通过 \texthref{https://msdn.microsoft.com/en-us/library/bb524831(v=vs.85}).aspx 下载安装。

\mySmallsection{0.3.5.3 线性代数模板库}

一些示例需要线性代数模板库，版本为3.3或更高。如果包管理器不提供\texttt{Eigen}，可以使用在线打包源(http://eigen.tuxfamily.org )安装它。例如，在GNU/Linux和macOS上，可以将\texttt{Eigen}安装到\texttt{\$HOME/Deps/Eigen}目录:

\begin{codebox}[shell]
$ eigen_version="3.3.4"
$ mkdir -p eigen
$ curl -Ls http://bitbucket.org/eigen/eigen/get/${eigen_version}.tar.gz | tar -xz -C eigen --strip-components=1
$ cd eigen
$ cmake -H. -Bbuild_eigen -
DCMAKE_INSTALL_PREFIX="$HOME/Deps/eigen" &> /dev/null
$ cmake --build build_eigen -- install &> /dev/null
\end{codebox}

\mySmallsection{0.3.5.4 Boost库}

\texttt{Boost}库适用于各种操作系统，大多数Linux发行版都通过它们的包管理器提供该库的安装。例如，在Ubuntu 14.04 LTS上，\texttt{Boost}文件系统库、\texttt{Boost Python}库和\texttt{Boost}测试库可以通过以下命令安装：

\begin{codebox}[shell]
$ sudo apt-get install libboost-filesystem-dev libboost-python-dev libboost-test-dev
\end{codebox}

对于macOS, \texttt{MacPorts}和自制程序都为最新版本的\texttt{Boost}提供了安装包。我们在macOS上的测试设置安装\texttt{Boost}如下：

\begin{codebox}[shell]
$ brew cask uninstall --force oclint
$ brew uninstall --force --ignore-dependencies boost
$ brew install boost
$ brew install boost-python3
\end{codebox}

Windows的二进制发行版也可以从\texttt{Boost}网站 \texthref{http://www.boost.org} 下载。或者，可以从 \texthref{https://www.boost.org} 下载源代码，并自己编译\texttt{Boost}库。

\mySmallsection{0.3.5.5 交叉编译器}

在类Debian/Ubuntu系统上，可以使用以下命令安装交叉编译器：

\begin{codebox}[shell]
$ sudo apt-get install gcc-mingw-w64 g++-mingw-w64 gfortran-mingw-w64
\end{codebox}

在macOS上，使用\texttt{Brew}，可以安装以下交叉编译器：

\begin{codebox}[shell]
$ brew install mingw-w64
\end{codebox}

其他包管理器提供相应的包。使用打包的跨编译器的另一种方法，是使用M交叉环境( \texthref{https://mxe.cc} )，并从源代码对其进行构建。

\mySmallsection{0.3.5.6 ZeroMQ, pkg-config, UUID和Doxygen}

Ubuntu 14.04 LTS上，这些包可以安装如下：

\begin{codebox}[shell]
$ sudo apt-get install pkg-config libzmq3-dev doxygen graphviz-dev uuid-dev
\end{codebox}

macOS上，我们建议使用\texttt{Brew}安装：

\begin{codebox}[shell]
$ brew install ossp-uuid pkg-config zeromq doxygen
\end{codebox}

\texttt{pkg-config}程序和\texttt{UUID}库只在类Unix系统上可用。 Windows上使用MSYS2环境，可以按如下方式安装这些依赖项(假设是64位版本)：

\begin{codebox}[shell]
$ pacman -S mingw64/mingw-w64-x86_64-zeromq
$ pacman -S mingw64/mingw-w64-x86_64-pkg-config
$ pacman -S mingw64/mingw-w64-x86_64-doxygen
$ pacman -S mingw64/mingw-w64-x86_64-graphviz
\end{codebox}

\mySmallsection{0.3.5.7 Conda的构建和部署}

想要使用\texttt{Conda}打包的示例的话，需要\texttt{Miniconda}和\texttt{Conda}构建和部署工具。\texttt{Miniconda}的安装说明之前已经给出。要在GNU/Linux和macOS上安装\texttt{Conda}构建和部署工具，请运行以下命令:

\begin{codebox}[shell]
$ conda install --yes --quiet conda-build anaconda-client jinja2 setuptools
$ conda clean -tipsy
$ conda info -a
\end{codebox}

这些工具也可以安装在Windows上:

\begin{codebox}[shell]
$conda_path = "C:\Deps\conda\Scripts\conda.exe"
$args = "install --yes --quiet conda-build anaconda-client jinja2 setuptools"
Start-Process -FilePath "$conda_path" -ArgumentList $args -Wait -Passthru
$args = "clean -tipsy"
Start-Process -FilePath "$conda_path" -ArgumentList $args -Wait -Passthru
$args = "info -a"
Start-Process -FilePath "$conda_path" -ArgumentList $args -Wait -Passthru
\end{codebox}


% --- recipe 4: 0.4-chinese.md ---
\section{0.4 测试环境}

示例在下列持续集成(CI)上进行过测试：

\begin{itemize}
\item Travis( \texthref{https://travis-ci.org} )用于GNU/Linux和macOS

\item Appveyor( \texthref{https://www.appveyor.com} )用于Windows

\item CircleCI ( \texthref{https://circleci.com} )用于附加的GNU/Linux测试和商业编译器

\end{itemize}

CI服务的配置文件可以在示例库中找到( \texthref{https://github.com/dev-cafe/cmake-cookbook/} ):

\begin{itemize}
\item Travis的配置文件为\texttt{travis.yml}

\item Appveyor的配置文件为\texttt{.appveyor.yml}

\item CircleCI的配置文件为\texttt{.circleci/config.yml}

\item Travis和Appveyor的其他安装脚本，可以在\texttt{ testing/dependencies}文件夹中找到。

\end{itemize}

\begin{myInfoBox}
\textbf{NOTE}：GNU/Linux系统上，Travis使用CMake 3.5.2和CMake 3.12.1对实例进行测试。macOS系统上用CMake 3.12.1进行测试。Appveyor使用CMake 3.11.3进行测试。Circle使用CMake 3.12.1进行测试。
\end{myInfoBox}

测试机制是一组Python脚本，包含在\texttt{testing}文件夹中。脚本\texttt{collect\_tests.py}将运行测试并报告它们的状态。示例也可以单独测试，也可以批量测试；\texttt{ collect\_tests.py}接受正则表达式作为命令行输入，例如:

\begin{codebox}[shell]
$ pipenv run python testing/collect_tests.py 'chapter-0[1,7]/recipe-0[1,2,5]'
\end{codebox}

该命令将对第1章和第7章的示例1、2和5进行测试。输出的示例如下:

\myGraphic{0.9}{images/chapter0/0-1.png}{}

要获得更详细的输出，可以设置环境变量\texttt{VERBOSE\_OUTPUT=ON}：

\begin{codebox}[shell]
$ env VERBOSE_OUTPUT=ON pipenv run python testing/collect_tests.py 'chapter-*/recipe-*'
\end{codebox}


% --- recipe 5: 0.5-chinese.md ---
\section{0.5 上报问题并提出改进建议}

请将遇到的问题反馈到 \texthref{https://github.com/dev-cafe/cmake-cookbook/issues。}

要对源码库进行贡献，我们建议对原始库 \texthref{https://github.com/dev-cafe/cmake-cookbook} 进行Fork，并使用Pull Request提交更改，可以参考这个页面 \texthref{https://help.github.com/articles/creating-a-pull-request-from-a-fork/} 。

对于非重要更改，我们建议在发送Pull Request之前，首先在 \texthref{https://github.com/devcafe/cmake-cookbook/issues} 上创建一个问题进行描述，并讨论所要更改的问题。


